\subsection{単冪Jordan型\texorpdfstring{$(3,1)$}{(3,1)}の全排除}
\label{subsec:cubic-jordan31-unipotent-vertex}
\begin{theorem}\label{thm:cubic-jordan31-unipotent-empty}
強い内在的H3の二次包絡の中山作用は、単一固有値1のJordan型 $(3,1)$ を持たない。
\end{theorem}
\begin{proof}
前節と同じねじり解除で $B=C^H$、$H=\exp N$、$C$ はCYとし、
\[
 Nx_1=x_0,\quad Nx_2=x_1,\quad Nx_0=Nx_3=0
\]
とする。元反巡回topの $N_{\rm tot}$ 核は18次元である。
次の18語の係数を順に $a_0,\ldots,a_{17}$ として一意に表示できる：
\[
\begin{gathered}
0001,0002,0003,0011,0012,0013,0022,0023,0031,\\
0032,0033,0113,0123,0133,0233,0303,0323,0333.
\end{gathered}
\]
middleの列22は $a_6x_0^2$ だけなので $a_6=0$。
$a_{12}=0$ なら列12が $a_4x_0^2$ だけとなり、$a_4=0$ も必要。
このとき $x_0,x_3$ を殺した全関係は $kx_1^2$ に含まれ、
$k\langle x_1,x_2\rangle/(x_1^2)$ への全射が指数成長を与える。
従って $a_{12}\ne0$、topを尺度変更して $a_{12}=1$ とする。

weightを $(0,1,2,0)$ とする。元六関係のunit minorは $-1$ であり、
最高weight部分は、可換三変数環 $k[x_0,x_1,x_2]$ 上のOre拡大
\[
 \operatorname{gr} C=k[x_0,x_1,x_2][x_3;\delta]
\]
の六関係となる。ただし $h=a_4$ として
\[
 \delta(x_0)=hx_0^2,\qquad \delta(x_1)=hx_0x_1,\qquad
 \delta(x_2)=-3hx_0x_2+2hx_1^2.
\]
残りの項はweightが厳密に低い。
Oreの並べ替えと低weightへの帰納で順序単項式が張り、AS4 Hilbert値と個数が一致するため基底となる。
従って上の随伴環同定は全次数で正しい。$h=0$ も含めている。
$N\delta=\delta N$、$Nx_3=0$ なので、PBW基底 $x_0^ix_1^jx_2^kx_3^l$ 上の
随伴 $N$ は可換三変数側だけに作用する。
その全三次核は
\[
 x_0^3,\ x_0^2x_3,\ x_0x_3^2,\ x_3^3,\quad
 x_0(2x_0x_2-x_1^2),\ (2x_0x_2-x_1^2)x_3
\]
で張られ、最高weightは2である。
輸送された正常三次 $F$ は $N(F)=0$ なので、先頭項の議論からweightが高々2となる。
従って $J=(x_0,x_1,x_3)$ に対し $F\in J^2$。
$H$ は $J$ を保ち、$B/J=k[x_2]$、従って $R=B/(F)$ に
$P=\pi k[x_2]$ が存在する。

前節の全tails自己拡大の議論をこの $P$ に適用する。
$x_2$ の作用を恒等にし、他の三生成元の非対角成分を
$f=(f_0,f_1,f_3)^t$ とする。
$J$ 内の三つの二次関係と $F\in J^2$ は一次化で零であり、残る三関係は可逆な漸化式を与える。
元18係数のまま、shiftの行列は
\[
 M=\begin{pmatrix}
 1&0&0\\
 \lambda&1&a_{16}-a_{14}\\
 3a_4&0&1
 \end{pmatrix},\qquad
 \lambda=1-2a_{11}+3a_{16}a_4-a_7-2a_9
\]
となる。特に $(M-I)^3=0$、$\det M=1$。
任意の初期値が一つの拡大を与えるので $\Ext^1(P,P)=k^3$。
$S=(-1)[2]$ と $P\simeq P(-1)$ に対し、
補題\ref{lem:cubic-odd-Serre-determinant}の行列式条件に矛盾する。
従って全係数・全退化境界を含めてこのJordan型はH3を与えない。
\end{proof}
公開再現器は\\
\texttt{verify\_cubic\_jordan31\_unipotent\_vertex.py}。
元18核、二つのdomain/growth列、unit source minor、最高weightのOre関係、
$N\delta=\delta N$、全三次核、twist後の全point拡大とmonodromyを検証する。
必要なSerre・tails論証は前節と上の通常証明に依存する。
単一固有値に残る非自明Jordan型は $(2,2),(2,1,1)$ であり、CY全被覆も未完成である。
